\begin{corollary}[7.4.3]
\label{IV.2.s7.4.3-fr}
\begin{enumerate}
\item[\emph{(i)}] Sea $A$ un $\boldsymbol P$-anillo local
noetheriano, y sea $A'=\widehat A$ su completado. Supongamos que,
para todo ideal primo $\mathfrak p'$ de $A'$, las fibras formales de
$A'_{\mathfrak p'}$ son geométricamente regulares; entonces, para
todo ideal primo $\mathfrak p$ de $A$, $A_{\mathfrak p}$ es un
$\boldsymbol P$-anillo.
\item[\emph{(ii)}] Supongamos que $\boldsymbol P$ satisface la
condición $(\mathrm P_{\mathrm{IV}})$ \textup{(7.3.6)}. Sean $A$ un
$\boldsymbol P$-anillo local noetheriano y $B$ una $A$-álgebra
local esencialmente de tipo finito, tal que el homomorfismo
$A\to B$ sea local. Pongamos $A'=\widehat A$, y sea
$\mathfrak n'$ el único ideal primo de $B'=B\otimes_A A'$ situado
sobre el ideal maximal de $B$ y el ideal maximal de $A'$. Si las
fibras formales de $B'_{\mathfrak n'}$ son geométricamente regulares,
entonces $B$ es un $\boldsymbol P$-anillo.
\end{enumerate}
\end{corollary}

\begin{proof}
\begin{enumerate}
\item[(i)] Como, por hipótesis,
$\operatorname{Spec}(A')\to\operatorname{Spec}(A)$ es un
$\boldsymbol P$-morfismo, también lo es
$\operatorname{Spec}(A'_{\mathfrak p'})\to
\operatorname{Spec}(A_{\mathfrak p})$ en virtud de (7.3.2, \emph{b)}),
para todo ideal primo $\mathfrak p$ de $A$ y todo ideal primo
$\mathfrak p'$ de $A'$ situado sobre $\mathfrak p$. Basta, pues,
aplicar el lema (7.4.2) a este morfismo, teniendo en cuenta que
$\operatorname{Spec}(A')\to\operatorname{Spec}(A)$ es sobreyectivo.
\item[(ii)] En virtud de (7.4.2), basta probar que el
morfismo
$\operatorname{Spec}(B'_{\mathfrak n'})\to\operatorname{Spec}(B)$
es un $\boldsymbol P$-morfismo. Ahora bien, se tiene
$B=C_{\mathfrak n}$, donde $C$ es una $A$-álgebra de tipo finito y
$\mathfrak n$ un ideal primo de $C$ situado sobre el ideal maximal
de $A$. Si se pone $C'=C\otimes_A A'$, de las hipótesis y de
(7.3.7) se deduce que
$\operatorname{Spec}(C')\to\operatorname{Spec}(C)$ es un
$\boldsymbol P$-morfismo; como $B'_{\mathfrak n'}$ es un anillo
local en un ideal primo de $C'$ situado sobre $\mathfrak n$, la
conclusión se deduce de (7.3.2, \emph{b)}).
\end{enumerate}
\end{proof}
